\documentclass{article}
\usepackage{amsmath,amssymb,hyperref}
\begin{document}
\section*{Determinants of prime alternating achiral knots}
\subsection{Definitions and mathematical statement}
For a smooth knot $K\subset S^3$, define
$\det(K)=|\Delta_K(-1)|$, where $\Delta_K$ is its Alexander polynomial.
This value is independent of multiplying the polynomial by $\pm t^j$.
A knot is prime if it is nontrivial and has no decomposition as a connected
sum of two nontrivial knots. It is alternating if it admits a diagram
whose crossings alternate over and under along the knot. It is achiral
if the unoriented knot is ambient isotopic to its mirror image.

Let
\[
 A=\{n\in\mathbb Z_{>0}:n\text{ is odd},\ n=a^2+b^2
 \text{ for some }a,b\in\mathbb Z\}\setminus\{1,9,49\}.
\]
Is every $n\in A$ equal to $\det(K)$ for some knot $K$ which is
simultaneously prime, alternating and achiral? Zero is allowed as one
of $a,b$; coprimality is not assumed.
\subsection{Short English statement}
Do prime alternating achiral knots realize every odd sum of two squares other than 1, 9 and 49 as a determinant?
\subsection{Sources}
A. Stoimenow, Problem 12.25 in T. Ohtsuki (ed.),
\emph{Problems on invariants of knots and 3-manifolds}, Geometry and Topology
Monographs 4 (2002), p. 541,
\url{https://msp.org/gtm/2002/04/gtm-2002-04-024s.pdf}.
See also A. Stoimenow, \emph{Square numbers, spanning trees and invariants
of achiral knots}, Communications in Analysis and Geometry 13 (2005),
591--631, \url{https://arxiv.org/abs/math/0003172}.
\end{document}
